% Options for packages loaded elsewhere
\PassOptionsToPackage{unicode}{hyperref}
\PassOptionsToPackage{hyphens}{url}
\PassOptionsToPackage{dvipsnames,svgnames,x11names}{xcolor}
%
\documentclass[
  11pt,
  a4paper,
]{ltjsarticle}
\usepackage{amsmath,amssymb}
\usepackage{iftex}
\ifPDFTeX
  \usepackage[T1]{fontenc}
  \usepackage[utf8]{inputenc}
  \usepackage{textcomp} % provide euro and other symbols
\else % if luatex or xetex
  \usepackage{unicode-math} % this also loads fontspec
  \defaultfontfeatures{Scale=MatchLowercase}
  \defaultfontfeatures[\rmfamily]{Ligatures=TeX,Scale=1}
\fi
\usepackage{lmodern}
\ifPDFTeX\else
  % xetex/luatex font selection
\fi
% Use upquote if available, for straight quotes in verbatim environments
\IfFileExists{upquote.sty}{\usepackage{upquote}}{}
\IfFileExists{microtype.sty}{% use microtype if available
  \usepackage[]{microtype}
  \UseMicrotypeSet[protrusion]{basicmath} % disable protrusion for tt fonts
}{}
\makeatletter
\@ifundefined{KOMAClassName}{% if non-KOMA class
  \IfFileExists{parskip.sty}{%
    \usepackage{parskip}
  }{% else
    \setlength{\parindent}{0pt}
    \setlength{\parskip}{6pt plus 2pt minus 1pt}}
}{% if KOMA class
  \KOMAoptions{parskip=half}}
\makeatother
\usepackage{xcolor}
\usepackage[margin=24mm]{geometry}
\setlength{\emergencystretch}{3em} % prevent overfull lines
\providecommand{\tightlist}{%
  \setlength{\itemsep}{0pt}\setlength{\parskip}{0pt}}
\setcounter{secnumdepth}{5}
\ifLuaTeX
  \usepackage{selnolig}  % disable illegal ligatures
\fi
\IfFileExists{bookmark.sty}{\usepackage{bookmark}}{\usepackage{hyperref}}
\IfFileExists{xurl.sty}{\usepackage{xurl}}{} % add URL line breaks if available
\urlstyle{same}
\hypersetup{
  colorlinks=true,
  linkcolor={blue},
  filecolor={Maroon},
  citecolor={Blue},
  urlcolor={blue},
  pdfcreator={LaTeX via pandoc}}

\author{}
\date{}

\begin{document}

\hypertarget{ux884cux5217}{%
\section{行列}\label{ux884cux5217}}

行列は数を長方形に並べたものですが、線形代数では
\textbf{線形写像の座標表現} として理解することが最重要です。

\[
A\in\mathbb{R}^{m\times n}
\]

は \texttt{R\^{}n} の入力を \texttt{R\^{}m} の出力へ写します。

\[
x\mapsto Ax
\]

\hypertarget{ux3064ux306eux898bux65b9}{%
\subsection{3つの見方}\label{ux3064ux306eux898bux65b9}}

\hypertarget{ux884cux306bux3088ux308bux898bux65b9}{%
\subsubsection{行による見方}\label{ux884cux306bux3088ux308bux898bux65b9}}

各出力成分は、入力との内積です。

\hypertarget{ux5217ux306bux3088ux308bux898bux65b9}{%
\subsubsection{列による見方}\label{ux5217ux306bux3088ux308bux898bux65b9}}

\texttt{Ax} は \texttt{A} の列の一次結合です。

\hypertarget{ux5199ux50cfux306bux3088ux308bux898bux65b9}{%
\subsubsection{写像による見方}\label{ux5199ux50cfux306bux3088ux308bux898bux65b9}}

空間全体を伸縮・せん断・回転・射影します。

この3つを行き来できることが重要です。

\end{document}
